\section{Future Work}
\label{sec:future}

The improvements below follow from the limitations in Section~\ref{subsec:limitations} or from a measured result. We consider the first two, Indian training data and an ensemble of the three models, the most useful for the effort they need.

\subsection{Data and the Indian Domain}
\label{subsec:fw-data}

The main next step is Indian training data with one box per item, either a relabeled part of DWSD or a new collection, with a part held out for testing. This addresses both the DWSD failure and the labeling mismatch behind part of it. Because the real-world data of \gentwo cost studio accuracy (Section~\ref{subsec:deployment}), new data should be added with the domains in mind, for example by keeping all studio images or by using a larger model. The weak OTHER class needs more examples, a split into clearer classes, or a merge into existing ones. Augmentations aimed at clutter, such as copy-paste~\cite{ghiasi2021simple}, smaller object scales, blur and lighting changes, would target the small, overlapping and half-hidden items of real-world photographs. The classes could also be arranged to follow the three streams of the SWM Rules 2016 (bio-degradable, non-bio-degradable and domestic hazardous)~\cite{moefcc2016solid}, so that each detection maps to a bin.

\subsection{Combining the Three Models}
\label{subsec:fw-ensemble}

No arm has the best point estimate on every \gentwo test set (Table~\ref{tab:gen2-test}). On real-world photographs the three \mapfifty values are within 0.012 of each other, but the arms differ more in the details: precision at the maximum-F1 confidence ranges from 0.366 (B) to 0.556 (C), PAPER AP@50 from 0.239 (A) to 0.330 (C), and random search is best on GLASS (0.296). These values have no intervals and we did not measure how often the models agree, but they suggest that the models make partly different mistakes. If so, combining them could beat each one alone without any extra training.

We would compare three ways of combining them. Box fusion pools the boxes of A, B and C and merges overlapping ones by weighted boxes fusion~\cite{solovyev2021weighted} or non-maximum suppression. Majority voting keeps a box only if at least two models find it (IoU $\ge 0.5$), which should raise precision; boxes found by one model could be flagged as uncertain instead of discarded. Weighted voting weights each model by its validation \mapfifty (not its test score, which would leak the test split). A search also leaves behind a whole population of good settings, the final swarm and the best random samples, so one could train the top $k$ of them and combine those. The converged swarm may not be diverse enough for this, so some of the best random samples should be mixed in. Each ensemble should be compared with the best single arm, chosen on validation, using the same paired bootstrap on every test set. The cost is inference time, which is about three times as long: roughly 100~ms per image on the laptop CPU against 31 to 38~ms for one \genone model.

\subsection{The Optimization Study}
\label{subsec:fw-opt}

Three training seeds per arm would let the intervals cover training noise, and several search seeds per method would show whether the PSO result repeats. The small proxy gaps among the top settings disappeared after full training, so a multi-fidelity search such as successive halving or Hyperband~\cite{li2018hyperband} could promote the best candidates to longer runs before a winner is chosen. The swarm converged early (in \gentwo it never improved after round~1), and PSO variants could keep it exploring for longer: the constriction factor of Clerc and Kennedy~\cite{clerc2002particle}, a ring topology, adaptive inertia, or simply more particles and rounds.

Other optimizers could be added as further arms against the same baselines at 32 evaluations: a genetic or evolutionary algorithm~\cite{young2015optimizing}, differential evolution~\cite{storn1997differential}, and Bayesian optimization such as the tree-structured Parzen estimator in Optuna~\cite{akiba2019optuna}. The protocol does not depend on PSO, and the beluga whale optimizer of Sayed et al.~\cite{sayed2024intelligent} could be tested in the same way. A wider search space (color augmentation, translation, loss gains, image size) would show whether six settings left too little room to beat the defaults, and a fitness scored on real-world validation images would aim the search at the photographs where tuning did not help. The proxy check should be repeated for \gentwo against a 100-epoch reference, together with a proxy run of the defaults to test the explanation in Section~\ref{subsec:why}. The statistics could be strengthened with Holm's correction, intervals for per-class AP and stored bootstrap resamples.

\subsection{Model and Deployment}
\label{subsec:fw-model}

Larger or newer detectors (YOLOv8s, YOLOv8m, YOLO11~\cite{jocher2024yolo11}) trained with the same settings would show how much of the real-world gap is due to model capacity, at some cost in speed on edge devices. Sliced inference~\cite{akyon2022slicing} on overlapping tiles suits wide photographs with many small items. Confidence calibration such as temperature scaling~\cite{guo2017calibration} would make a score of 0.5 mean roughly 50\% correct, which matters when a bin acts on a threshold. Per-class confidence thresholds chosen on validation data could replace the single threshold, since real-world precision is low and AP varies widely by class.

Our latencies come from a laptop CPU and a cloud GPU only. The models should be exported to ONNX or TFLite and timed on a Raspberry Pi, the kind of hardware that Alshammari et al.~\cite{alshammari2025smart} use for a smart-bin classifier, and tested on live video, since real sorting is continuous.
